% ============================================================
\section{Ejercicio 1}

% ============================================================

\begin{tcolorbox}[enunciado]
Considera el siguiente par de problemas:
\begin{itemize}
    \item \textbf{Problema:} 3SAT-IGUALDAD. \\ 
    \textbf{Ejemplar:} Una fórmula booleana $\phi$ en FNC en la cual el número de cláusulas es igual al número de variables. \\
    \textbf{Pregunta:} ¿Existe una asignación que satisface $\phi$?
    \item \textbf{Problema:} 3SAT BALANCEADO. \\
    \textbf{Ejemplar:} Una fórmula booleana $\phi$ en FNC en la cual cada cláusula contiene 3 literales, y cada variable en $\phi$ aparece negada y sin negación el mismo número de veces. \\
    \textbf{Pregunta:} ¿Existe una asignación que satisface $\phi$?
\end{itemize}
\begin{enumerate}[label=(\alph*)]
    \item Utilizando reducciones, demuestra la equivalencia de los problemas. En cada caso, demuestra que se cumple la reducción correspondiente y ejemplifica las transformaciones correspondientes. Indica qué técnicas de reducción se utilizan en las demostraciones.
    \item Demuestra que ambos problemas son $NP-Completos.$
\end{enumerate}
\end{tcolorbox}

% ------------------------- Solución -------------------------
\section{Respuestas}
\subsection{a}

\subsection{b}